\documentclass{beamer}
\usetheme{Madrid}  % or Berlin, AnnArbor, etc.
\usecolortheme{crane}
\usepackage{booktabs} % better tables


\usepackage{colortbl}
\definecolor{chip}{RGB}{255, 220, 180}
\definecolor{sum}{RGB}{180, 220, 255}
\definecolor{prod}{RGB}{180, 255, 200}
\definecolor{sbc}{RGB}{240, 180, 255}

\title{Expansion of Aggregation Dynamics with Master Equations}
\author{Fitz Koch}
\date{\today}

\begin{document}
\begin{frame}
    \titlepage
\end{frame}

\begin{frame}{Outline}
    \tableofcontents
\end{frame}

\section{Purpose}
\begin{frame}{Purpose}
    \begin{itemize}
        \item Key Question: How do we avoid gelation in aggregation dynamics with the \textit{least costly} regulatory interference?
        \item Imagine corporate consolidation, or anywhere else where things tend to get big and we don't want them to. How do we prevent an entity from getting \textit{overwhelmingly} big, ala the giant component, with the \textit{least} possible expenditure of effort?
        \item Cost:  \textbf{number of interventions}, \textbf{size of interventions}
    \end{itemize}
\end{frame}

\section{Setup}
\begin{frame}{Aggregation}
    \begin{columns}
        \begin{column}{0.5\textwidth}
            Aggregation is generally the transition:
            $$
                A_{i} + A_{j} \xrightarrow[]{K_{ij}} A_{i+j}
            $$
            Yields master equation:
            \begin{align*}
                \frac{\mathrm{d} c_{k}}{\mathrm{d} t }
                 & = \frac{1}{2} \sum_{i+j=k} K_{ij} c_{i} c_{j} - c_{k} \sum_{i \geq  1} K_{ik}c_{i}\,                            \\
                 & = \frac{1}{2} \sum_{i, j \leq  1} K_{ij} c_{i} c_{j} \Big[\delta_{i+j, k} - \delta_{i, k} - \delta_{j, k} \Big]
            \end{align*}
        \end{column}
        \begin{column}{0.5\textwidth}
            Definitions:
            \begin{itemize}
                \item $A$ is a cluster
                \item $i$ and $j$ are masses
                \item $K_{ij}$ is the rate of aggregation between clusters of mass $i$ and mass $j$ into a clusters of mass $i +j$
                \item $K$, then, is an infinite symmetric matrix $K_{ij}=K_{ji}$
                \item $\delta$ is the Kronecker Delta
            \end{itemize}
        \end{column}
    \end{columns}
\end{frame}

\begin{frame}{Fragmentation}
    \begin{columns}
        \begin{column}{0.5\textwidth}
            Fragmentation is generally the transition:
            $$
                A_{i+j}  \xrightarrow[]{F_{ij}} A_{i} + A_{j}
            $$
            Yields master equation:
            \begin{align*}
                \frac{\mathrm{d} c_{k}}{\mathrm{d} t } = \sum _{j \geq  1} F_{kj}c_{j+k} - \frac{1}{2} \sum _{i+j=k}F_{ij}
            \end{align*}
        \end{column}
        \begin{column}{0.5\textwidth}
            Definitions:
            \begin{itemize}
                \item $A$ is a cluster
                \item $i$ and $j$ are masses
                \item $F_{ij}$ is the rate of fragmentation of a cluster of size $i+j$ into a cluster of size $i$ and another of size $j$
            \end{itemize}
        \end{column}
    \end{columns}
\end{frame}


\begin{frame}{Combined Master Equations}
    Put together, we have the combined master equation:
    $$
        \frac{\mathrm{d} c_{k}}{\mathrm{d} t }  = \frac{1}{2} \sum _{i+j=k} K_{ij}c_{i}c_{j} - c_{k}\sum _{j \geq  1}K_{jk}c_{j} + \sum _{j\geq  1} F_{kj} c_{j+k} - \frac{1}{2} c_{k} \sum _{i+j=k}F_{ij}
    $$
    With a competition between the aggregation kernel $ K $  and the fragmentation kernel $ F $.

    We can now be more specific about our two definitions of the `cost' of fragmentation. Let $ l_i $, the cost per intervention, be the number of times that the aggregation term $ F $  applies. Let $ l_s $ instead be the sum of $ \min(i, j) $ over all the timesteps $ F $ is applied.
\end{frame}

\begin{frame}{Choosing Aggregation Kernel}
    \small
    $$
        \frac{\mathrm{d} c_{k}}{\mathrm{d} t} = \frac{1}{2} \sum_{i+j=k} K_{ij}c_{i}c_{j} - c_{k}\sum_{j \geq 1}K_{jk}c_{j} + \sum_{j\geq 1} F_{kj} c_{j+k} - \frac{1}{2} c_{k} \sum_{i+j=k}F_{ij}
    $$
    \normalsize
    Based on polymerization, parameterize aggregation on number of bond points $ f $ :
    $$
        K_{ij} = (f-2)^{2}ij +2(f-2)(i+j) + 4
    $$
    \begin{itemize}
        \item Dunbar's number: $ f=152 \implies K_{ij} = \alpha(150^2ij + 300(i+j) + 4)$
        \item Proposed \textit{homo corporatum} number: $ f \to \infty \implies K_{ij}  \approx \alpha(ij) $
        \item Note: $ K_{ij} = ij $ is the same as percolation // erdos-renyi graph.
    \end{itemize}
\end{frame}


\begin{frame}{Choosing Fragmentation Kernel }
    \small
    $$
        \frac{\mathrm{d} c_{k}}{\mathrm{d} t} = \frac{1}{2} \sum_{i+j=k} K_{ij}c_{i}c_{j} - c_{k}\sum_{j \geq 1}K_{jk}c_{j} + \sum_{j\geq 1} F_{kj} c_{j+k} - \frac{1}{2} c_{k} \sum_{i+j=k}F_{ij}
    $$
    \normalsize

    Parameterized by $\gamma$, a few different models:

    \vspace{0.5em}
    \begin{tabular}{ll}
        \toprule
        \textbf{Name}                        & \textbf{Equation}                                               \\
        \midrule
        \rowcolor{chip}  Chipping Model      & $F_{ij} = \gamma_{0} + \gamma_{1}(\delta_{i,1} + \delta_{j,1})$ \\
        \rowcolor{sum}   Sum Size            & $F_{ij} = \gamma(i + j)$                                        \\
        \rowcolor{prod}  Product Splitting   & $F_{ij} = \gamma(ij)$                                           \\
        \rowcolor{sbc}   Size Based Chipping & $F_{ij} = \gamma(j\delta_{i,1} + i\delta_{j,1})$                \\
        \bottomrule
    \end{tabular}
    \vspace{0.5em}

    Connect back to goal: how well does each prevent gelation? How much does the intervention of each cost?

\end{frame}

\section{Simulation}

\end{document}